%%%%%
%
%% Dans ce fichier, les �ventuels accents sont cod�s
%% en ISO-8859-15 (aka latin9)
%% La ligne suivante donne trois exemples 
%% e accent aigu : �, e accent grave : �, e-dans-l'o : �
\section{Correction de la hauteur GPS par rapport a l'ellipsoide
 WGS84 avec l'ondulation du geoide}

\itemstyle{Op�rateur~:} \operfmt{Cor\_ondulation\_geoide84}
\index{Cor\_ondulation\_geoide84}

%%%%%OBJET%%%%
\itemstyle{Objet~:}
Corrige la hauteur GPS par rapport � l'ellipsoide WGS84
 avec l'ondulation du geoide pour le lieu.

%%%%%ARGUMENTS%%%%
\bigskip
%% Dans ce fichier, les �ventuels accents sont cod�s
%% en ISO-8859-15 (aka latin9)
%% La ligne suivante donne trois exemples 
%% e accent aigu : �, e accent grave : �, e-dans-l'o : �
%%%%%ENTREES%%%%
\begin{tabarg}
\multicolumn{3}{l}{\titretab {Entr\'ees}}\\
\hline
$Lat$ & Vecteur, \texttt{float} & Latitude (degres)\\
$Lon$ & Vecteur, \texttt{float} & Longitude (degres)\\
$Alt$ & Vecteur, \texttt{float} & Altitude GPS (metres)\\
\end{tabarg}

\medskip
%%%%%SORTIES%%%%
\begin{tabarg}
\multicolumn{3}{l}{\titretab {Sorties}}\\
\hline
$Alt$ & Vecteur, \texttt{float} & Altitude corrig�e (metres)\\
\end{tabarg}


%%%%%COMMENTAIRE%%%%
\itemstyle{Commentaire~:}
Le fichier EGM180.NOR contient des coefficients g�opotentielles normalis�s 
(Cnm, SNM) pour le mod�le WGS 84. Il est utilis� comme fichier d'entr�e pour 
l'op�rateur.
Sauvegard� sous le r�pertoire aida/extern/EGM\_84, il doit �stre copi� sous
 aida/share/common/EGM180.NOR.

%%%%%AUTEUR%%%%
\itemstyle{Auteur~:}
Le programme \textsc{fortran} qui fait effectivement les calculs
a �t� charg� sur le site du National Geospatial-intelligence Agency (NGA).
L'enrobage sous forme d'op�rateur a �t� fait par P. Richard.

%%%%%REFERENCE%%%%
\itemstyle{R�f�rences~:}
\url{http://earth-info.nga.mil/GandG/wgs84/gravitymod/wgs84_180/wgs84_180.html}
